\chapter{Killing Fields}
\label{chap:pet-killing-fields}

Killing fields encode infinitesimal isometries. Their first-order equations control
the isometry algebra, while curvature assumptions yield vanishing, fixed-point,
connectivity, and symmetry-rank results.

\section{Killing Fields in General}

The metric Lie derivative and covariant derivative give equivalent local
descriptions of Killing fields.

\begin{definition}[Killing field]
  \label{def:pet-ch8-killing-field}
  \lean{PetersenLib.IsKillingField}
  \source{petersen:page-0327}
  A vector field $X$ on a Riemannian manifold $(M,g)$ is called a \emph{Killing
  field} if the local flows generated by $X$ act by isometries.
\end{definition}

\begin{proposition}[metric Lie-derivative criterion]
  \label{prop:pet-ch8-killing-lie-derivative}
  \lean{PetersenLib.isKillingField_iff_lieDerivative_metric_eq_zero}
  \source{petersen:page-0327,petersen:page-0328}
  \uses{def:pet-ch8-killing-field}
  \dcref{ch8:8.1.1}
  A vector field $X$ on a Riemannian manifold $(M,g)$ is a Killing field if and
  only if $L_Xg=0$.
\end{proposition}
\begin{proof}
  \uses{prop:pet-ch8-killing-lie-derivative}
  \sketch
  Let \(F^t\) be the local flow of \(X\). For \(v,w\in T_pM\),
  \[
    \frac{d}{dt}g(DF^t v,DF^t w)\big|_{t=t_0}
    =(L_Xg)(DF^{t_0}v,DF^{t_0}w).
  \]
  Thus \(L_Xg=0\) exactly when the displayed inner product is constant in
  \(t\). Its value at \(t=0\) is \(g(v,w)\), so this is equivalent to every
  \(F^t\) being a local isometry.
\end{proof}

\begin{proposition}[covariant-derivative criterion]
  \label{prop:pet-ch8-killing-skew-symmetric}
  \lean{PetersenLib.isKillingField_iff_isSkewAdjoint_covariantDerivative}
  \source{petersen:page-0328}
  \uses{def:pet-ch8-killing-field, prop:pet-ch8-killing-lie-derivative}
  \dcref{ch8:8.1.2}
  A vector field $X$ is a Killing field if and only if $v\mapsto\nabla_vX$ is a
  skew-symmetric $(1,1)$-tensor, i.e.\ $g(\nabla_vX,w)=-g(\nabla_wX,v)$ for all
  $v,w$.
\end{proposition}
\begin{proof}
  \uses{prop:pet-ch8-killing-skew-symmetric}
  \sketch
  Metric compatibility and torsion-freeness give
  \[
    (L_Xg)(V,W)=g(\nabla_VX,W)+g(V,\nabla_WX).
  \]
  Hence \cref{prop:pet-ch8-killing-lie-derivative} identifies the Killing
  equation with skew-adjointness of the endomorphism \(V\mapsto\nabla_VX\).
\end{proof}

\begin{proposition}[curvature identities for Killing fields]
  \label{prop:pet-ch8-killing-curvature-identities}
  \lean{PetersenLib.killingField_hessian_eq_curvature,PetersenLib.killingField_bracket_curvature}
  \source{petersen:page-0328,petersen:page-0329}
  \uses{def:pet-ch8-killing-field}
  \dcref{ch8:8.1.3}
  If $X\in\iso(M,g)$, then
  \[
    \nabla^2_{V,W}X=-R(X,V)W.
  \]
  If $X,Y\in\iso(M,g)$, then
  \[
    [\nabla X,\nabla Y](V)+\nabla_V[X,Y]=R(X,Y)V.
  \]
\end{proposition}
\begin{proof}
  \uses{prop:pet-ch8-killing-curvature-identities}
  \sketch
  An isometric flow preserves the Levi--Civita connection, so
  \(L_X\nabla=0\). The standard variation identity
  \[
    (L_X\nabla)_VW=R(X,V)W+\nabla^2_{V,W}X
  \]
  yields the first formula. Expanding \(\nabla_V[X,Y]\) and applying that
  formula once to \(X\) and once to \(Y\) gives
  \[
    [\nabla X,\nabla Y](V)+\nabla_V[X,Y]
    =-R(Y,V)X+R(X,V)Y=R(X,Y)V,
  \]
  where the last equality is the first Bianchi identity.
\end{proof}

\begin{proposition}[determination by the first jet]
  \label{prop:pet-ch8-killing-uniqueness}
  \lean{PetersenLib.killingField_eq_zero_of_jet_eq_zero,PetersenLib.isoExactSequence}
  \source{petersen:page-0329}
  \uses{def:pet-ch8-killing-field, prop:pet-ch8-killing-skew-symmetric}
  \dcref{ch8:8.1.4}
  For a given $p\in M$, a Killing field $X\in\iso(M,g)$ is uniquely determined by
  $X|_p$ and $(\nabla X)|_p$. In particular there is a short exact sequence
  \[
    0\to\iso_p\to\iso(M,g)\to t_p\to0,
  \]
  where $\iso_p=\{X\in\iso(M,g)\mid X|_p=0\}$ and
  $t_p=\{X|_p\in T_pM\mid X\in\iso(M,g)\}$.
\end{proposition}
\begin{proof}
  \uses{prop:pet-ch8-killing-uniqueness}
  \sketch
  Along every geodesic \(c\), an infinitesimal isometric variation makes
  \(X|_c\) a Jacobi field. Consequently \(X(p)=0\) and
  \((\nabla X)_p=0\) force \(X\) to vanish on a normal neighborhood of \(p\);
  continuation over the connected manifold gives \(X=0\). Therefore
  \(X\mapsto(X_p,(\nabla X)_p)\) is injective. The evaluation map
  \(\iso(M,g)\to T_pM\) has kernel \(\iso_p\) and image \(t_p\), which proves
  exactness of the displayed sequence.
\end{proof}

\begin{theorem}[geometry of the zero set]
  \label{thm:pet-ch8-zero-set-totally-geodesic}
  \lean{PetersenLib.killingField_zeroSet_totallyGeodesic_evenCodim}
  \source{petersen:page-0330}
  \uses{def:pet-ch8-killing-field, prop:pet-ch8-killing-skew-symmetric}
  \dcref{ch8:8.1.5}
  The zero set of a Killing field is a disjoint union of totally geodesic
  submanifolds, each of even codimension.
\end{theorem}
\begin{proof}
  \uses{thm:pet-ch8-zero-set-totally-geodesic}
  \sketch
  The zero set is the common fixed set of the local isometries generated by
  \(X\), hence each component is totally geodesic. At a zero \(p\), its tangent
  space is \(\ker(\nabla X)_p\). By
  \cref{prop:pet-ch8-killing-skew-symmetric}, \((\nabla X)_p\) is
  skew-adjoint, so its rank, equal to the component's codimension, is even.
\end{proof}

\begin{theorem}[dimension and integration of the isometry algebra]
  \label{thm:pet-ch8-isometry-lie-algebra-dimension}
  \lean{PetersenLib.isoLieAlgebra,PetersenLib.isoLieAlgebra_dim_le,PetersenLib.isoLieAlgebra_eq_lieAlgebra_isoGroup}
  \source{petersen:page-0330,petersen:page-0331}
  \uses{def:pet-ch8-killing-field, prop:pet-ch8-killing-uniqueness}
  \dcref{ch8:8.1.6}
  The set of Killing fields $\iso(M,g)$ is a Lie algebra of dimension
  $\le\binom{n+1}{2}$. Furthermore, if $M$ is complete, then $\iso(M,g)$ is the
  Lie algebra of $\mathrm{Iso}(M,g)$.
\end{theorem}
\begin{proof}
  \uses{thm:pet-ch8-isometry-lie-algebra-dimension,lem:pet-ch8-killing-flow-oneparameter-subgroup}
  \sketch
  The identity \(L_{[X,Y]}=[L_X,L_Y]\) makes \(\iso(M,g)\) a Lie algebra.
  By \cref{prop:pet-ch8-killing-uniqueness}, the first-jet map embeds it into
  \(T_pM\oplus\mathfrak{so}(T_pM)\); therefore
  \[
    \dim\iso(M,g)\le n+\frac{n(n-1)}2=\binom{n+1}{2}.
  \]
  On a complete manifold, \cref{lem:pet-ch8-killing-flow-oneparameter-subgroup}
  identifies these fields with the one-parameter subgroups of
  \(\mathrm{Iso}(M,g)\).
\end{proof}

\begin{lemma}[Lie group structure of the isometry group and the one-parameter-subgroup correspondence]
  \label{lem:pet-ch8-killing-flow-oneparameter-subgroup}
  \lean{PetersenLib.isoGroup_isLieGroup,PetersenLib.killingFlow_oneParameterSubgroup_correspondence}
  \source{petersen:page-0330,petersen:page-0331}
  \uses{def:pet-ch8-killing-field}
  Endow $\mathrm{Iso}(M,g)$ with the compact-open topology, so convergence means
  uniform convergence on compact sets. If $M$ is complete, this topology makes
  $\mathrm{Iso}(M,g)$ into a Lie group whose Lie algebra is $\iso(M,g)$: the
  flows of Killing fields are then defined for all time and consist of
  isometries, yielding differentiable one-parameter subgroups of
  $\mathrm{Iso}(M,g)$, and conversely every differentiable one-parameter
  subgroup of $\mathrm{Iso}(M,g)$ arises as the flow of a Killing field; this
  correspondence identifies the Lie algebra of $\mathrm{Iso}(M,g)$ with
  $\iso(M,g)$.
\end{lemma}
\begin{proof}
  \uses{lem:pet-ch8-killing-flow-oneparameter-subgroup}
  \sketch
  With the compact-open topology, the isometry group of a complete Riemannian
  manifold is a Lie group. Completeness makes every Killing flow global, and
  \(t\mapsto F^t\) is then a one-parameter subgroup of this group. Conversely,
  differentiating a one-parameter subgroup of isometries produces a vector
  field whose flow preserves \(g\). These two operations are inverse and
  identify the two Lie algebras.
\end{proof}

\begin{remark}[maximal-dimensional isometry groups among space-form quotients]
  \label{rem:pet-ch8-max-isometry-space-forms}
  \lean{PetersenLib.maxIsometryDim_spaceForm_quotient}
  \source{petersen:page-0331}
  \uses{thm:pet-ch8-isometry-lie-algebra-dimension}
  For the simply connected space forms $S^n_0$ (of constant curvature $0$),
  $\dim(\mathrm{Iso}(S^n_0))=\binom{n+1}2$, the maximal possible value. Among
  other complete constant-curvature manifolds $S^n_0/\Gamma$ with
  $\Gamma\subset\mathrm{Iso}(S^n_0)$ acting freely and discontinuously, the
  isometry group has maximal dimension $\binom{n+1}2$ only for
  $\Gamma=\{\pm I\}$ (the sphere itself is $\Gamma=\{I\}$), i.e.\ only for
  $\mathbb{RP}^n$ in addition to $S^n_0$ itself.
\end{remark}
\begin{proof}
  \uses{rem:pet-ch8-max-isometry-space-forms}
  \sketch
  Killing fields on \(S^n_0/\Gamma\) are the \(\Gamma\)-invariant fields on
  \(S^n_0\). Maximal dimension forces their connected group to be
  \(\mathrm{Iso}(S^n_0)_0\), so every element of \(\Gamma\) centralizes that
  group. The centralizer consists only of the identity and, in the spherical
  case, the antipodal map. Since the latter is the only nontrivial such element
  acting freely, the only additional quotient is \(\mathbb{RP}^n\).
\end{proof}

\begin{theorem}[maximal symmetry implies constant curvature]
  \label{thm:pet-ch8-maximal-symmetry-constant-curvature}
  \lean{PetersenLib.isoLieAlgebra_dim_eq_max_iff_constantCurvature}
  \source{petersen:page-0331}
  \uses{thm:pet-ch8-isometry-lie-algebra-dimension, prop:pet-ch8-killing-uniqueness}
  If $\dim(\iso(M,g))=\binom{n+1}2$, then $(M,g)$ has constant curvature.
\end{theorem}
\begin{proof}
  \uses{thm:pet-ch8-maximal-symmetry-constant-curvature}
  \sketch
  Equality in the dimension bound makes the first-jet injection from
  \cref{thm:pet-ch8-isometry-lie-algebra-dimension} an isomorphism
  \[
    \iso(M,g)\simeq T_pM\oplus\mathfrak{so}(T_pM).
  \]
  Fields with zero value at \(p\) realize every infinitesimal rotation, so the
  isotropy acts transitively on \(2\)-planes and all sectional curvatures at
  \(p\) agree. Fields whose prescribed jet is \((v,0)\) move \(p\) through a
  neighborhood, making that curvature value locally constant. Connectedness
  then makes it constant on \(M\).
\end{proof}

\section{Killing Fields in Negative Ricci Curvature}

The norm of a Killing field satisfies identities that turn nonpositive Ricci
curvature into rigidity through the maximum principle.

\begin{proposition}[gradient, Hessian, and Laplacian of a Killing-field norm]
  \label{prop:pet-ch8-killing-norm-formulas}
  \lean{PetersenLib.killingField_normSq_gradient,PetersenLib.killingField_normSq_hessian,PetersenLib.killingField_normSq_laplacian}
  \source{petersen:page-0332}
  \uses{def:pet-ch8-killing-field, prop:pet-ch8-killing-skew-symmetric, prop:pet-ch8-killing-curvature-identities}
  \dcref{ch8:8.2.1}
  Let $X$ be a Killing field on $(M,g)$ and $f=\tfrac12g(X,X)=\tfrac12|X|^2$.
  Then
  \begin{enumerate}
    \item $\nabla f=-\nabla_XX$.
    \item $\mathrm{Hess}\,f(V,V)=|\nabla_VX|^2-R(V,X,X,V)$.
    \item $\Delta f=|\nabla X|^2-\mathrm{Ric}(X,X)$.
  \end{enumerate}
\end{proposition}
\begin{proof}
  \uses{prop:pet-ch8-killing-norm-formulas}
  \sketch
  For \(f=\tfrac12|X|^2\), skew-adjointness from
  \cref{prop:pet-ch8-killing-skew-symmetric} gives
  \(D_Vf=g(\nabla_VX,X)=-g(V,\nabla_XX)\), proving the gradient identity.
  Differentiate once more and use
  \(\nabla^2_{V,W}X=-R(X,V)W\) to obtain
  \[
    \operatorname{Hess}f(V,V)=|\nabla_VX|^2-R(V,X,X,V).
  \]
  Tracing over an orthonormal frame yields
  \(\Delta f=|\nabla X|^2-\operatorname{Ric}(X,X)\).
\end{proof}

\begin{theorem}[Bochner vanishing for nonpositive Ricci curvature]
  \label{thm:pet-ch8-bochner-negative-ricci}
  \lean{PetersenLib.bochner_negativeRicci_killingParallel}
  \source{petersen:page-0333}
  \uses{prop:pet-ch8-killing-norm-formulas}
  \dcref{ch8:8.2.2}
  Suppose $(M,g)$ is compact with $\mathrm{Ric}\le0$. Then every Killing field
  is parallel. Furthermore, if $\mathrm{Ric}<0$, then there are no nontrivial
  Killing fields.
\end{theorem}
\begin{proof}
  \uses{thm:pet-ch8-bochner-negative-ricci}
  \sketch
  For a Killing field \(X\), \cref{prop:pet-ch8-killing-norm-formulas} gives
  \[
    \Delta\!\left(\tfrac12|X|^2\right)
    =|\nabla X|^2-\operatorname{Ric}(X,X)\ge0.
  \]
  Compactness and the maximum principle force this function and its Laplacian
  to be constant, hence \(\nabla X=0\). If \(\operatorname{Ric}<0\), the same
  equality also forces \(X=0\).
\end{proof}

\begin{corollary}[isometry-group bound under nonpositive Ricci curvature]
  \label{cor:pet-ch8-isometry-dimension-bound-negative-ricci}
  \lean{PetersenLib.isoDim_le_dim_negativeRicci,PetersenLib.isoGroup_finite_negativeRicci}
  \source{petersen:page-0333}
  \uses{thm:pet-ch8-bochner-negative-ricci}
  \dcref{ch8:8.2.3}
  With $(M,g)$ as in \cref{thm:pet-ch8-bochner-negative-ricci},
  \[
    \dim(\iso(M,g))=\dim(\mathrm{Iso}(M,g))\le\dim M,
  \]
  and $\mathrm{Iso}(M,g)$ is finite if $\mathrm{Ric}(M,g)<0$.
\end{corollary}
\begin{proof}
  \uses{cor:pet-ch8-isometry-dimension-bound-negative-ricci}
  \sketch
  By \cref{thm:pet-ch8-bochner-negative-ricci}, every Killing field is
  parallel, so evaluation at one point embeds \(\iso(M,g)\) into \(T_pM\).
  The identification in
  \cref{thm:pet-ch8-isometry-lie-algebra-dimension} gives the group-dimension
  statement. Under strict negativity the Lie algebra vanishes; the compact
  isometry group of compact \(M\) is therefore discrete and hence finite.
\end{proof}

\begin{corollary}[Euclidean splitting of the universal cover]
  \label{cor:pet-ch8-splitting-negative-ricci}
  \lean{PetersenLib.universalCover_splits_negativeRicci}
  \source{petersen:page-0333}
  \uses{thm:pet-ch8-bochner-negative-ricci}
  \dcref{ch8:8.2.4}
  With $(M,g)$ as in \cref{thm:pet-ch8-bochner-negative-ricci} and
  $k=\dim(\iso(M,g))$, the universal cover splits isometrically as
  $\widetilde M=\mathbb{R}^k\times N$.
\end{corollary}
\begin{proof}
  \uses{cor:pet-ch8-splitting-negative-ricci}
  \sketch
  The \(k\) Killing fields lift to linearly independent parallel fields on
  \(\widetilde M\). Their span is a parallel flat distribution and its
  orthogonal complement is parallel. The de Rham splitting theorem, applied
  on the complete simply connected cover, gives
  \(\widetilde M\cong\mathbb{R}^k\times N\) isometrically.
\end{proof}

\begin{corollary}[flatness of homogeneous metrics with nonpositive Ricci curvature]
  \label{cor:pet-ch8-homogeneous-negative-ricci-flat}
  \lean{PetersenLib.homogeneous_negativeRicci_flat}
  \source{petersen:page-0333}
  \uses{thm:pet-ch8-bochner-negative-ricci}
  \dcref{ch8:8.2.5}
  A compact homogeneous space with $\mathrm{Ric}\le0$ is flat. In particular,
  any Ricci-flat homogeneous space is flat.
\end{corollary}
\begin{proof}
  \uses{cor:pet-ch8-homogeneous-negative-ricci-flat}
  \sketch
  By \cref{thm:pet-ch8-bochner-negative-ricci}, all Killing fields are
  parallel. Homogeneity makes their values span every tangent space. Thus
  every tangent vector extends locally to a parallel field, which forces the
  curvature tensor to vanish.
\end{proof}

\begin{theorem}[quasi-negative Ricci curvature excludes Killing fields]
  \label{thm:pet-ch8-quasi-negative-ricci-no-killing}
  \lean{PetersenLib.quasiNegativeRicci_noNontrivialKillingField}
  \source{petersen:page-0334}
  \uses{thm:pet-ch8-bochner-negative-ricci, prop:pet-ch8-killing-norm-formulas}
  \dcref{ch8:8.2.6}
  Suppose $(M,g)$ is a compact manifold with \emph{quasi-negative} Ricci
  curvature, i.e.\ $\mathrm{Ric}\le0$ everywhere and $\mathrm{Ric}(v,v)<0$ for
  all $v\in T_pM\setminus\{0\}$ for some $p\in M$. Then $(M,g)$ admits no
  nontrivial Killing field.
\end{theorem}
\begin{proof}
  \uses{thm:pet-ch8-quasi-negative-ricci-no-killing}
  \sketch
  A Killing field is parallel by
  \cref{thm:pet-ch8-bochner-negative-ricci}. If it were nonzero, it would be
  nonzero at the point where Ricci curvature is negative definite. But
  \cref{prop:pet-ch8-killing-norm-formulas} would give there
  \(0=\Delta(\tfrac12|X|^2)=-\operatorname{Ric}(X,X)>0\), a contradiction.
\end{proof}

\begin{definition}[$F$-structure]
  \label{def:pet-ch8-f-structure}
  \lean{PetersenLib.FStructure}
  \source{petersen:page-0334}
  \uses{def:pet-ch8-killing-field}
  An \emph{$F$-structure} on $M$ consists of a finite covering of open sets
  $U_i$ of some finite covering space $\widetilde M\to M$, together with a
  Killing field $X_i$ on each $U_i$, such that whenever $U_i\cap U_j\ne\varnothing$
  the fields commute there: $[X_i,X_j]=0$ on $U_i\cap U_j$. The $F$-structure is
  \emph{pure} of positive rank when, roughly, the collection of local Killing
  fields is nowhere all simultaneously trivial in a suitable sense, measured via
  the function $f=\det[g(X_i,X_j)]$ (which reduces to $f=g(X,X)$ when only one
  field is given globally).
\end{definition}

\begin{remark}[Rong's generalization of Bochner's theorem]
  \label{rem:pet-ch8-rong-f-structure}
  \lean{PetersenLib.rong_pureFStructure_negativeRicci}
  \source{petersen:page-0334}
  \uses{def:pet-ch8-f-structure, thm:pet-ch8-bochner-negative-ricci}
  X. Rong generalized \cref{thm:pet-ch8-bochner-negative-ricci} to show that a
  closed Riemannian manifold with negative Ricci curvature admits no pure
  $F$-structure of positive rank. The strategy mirrors the proof of
  \cref{thm:pet-ch8-bochner-negative-ricci}: one shows the function
  $f=\det[g(X_i,X_j)]$ from \cref{def:pet-ch8-f-structure} is a well-defined,
  sufficiently regular function on $M$ admitting a Bochner-type formula.
\end{remark}

\section{Killing Fields in Positive Curvature}

In positive sectional curvature, zeros of Killing fields relate the topology
of fixed-point sets to the rank of commuting isometric actions.

\begin{remark}[the Hopf conjecture and low-dimensional cases]
  \label{rem:pet-ch8-hopf-conjecture-low-dimensions}
  \lean{PetersenLib.hopfConjecture_lowDimensions}
  \source{petersen:page-0334}
  Every vector field on an even-dimensional sphere has a zero, since
  $\chi(S^{2m})=2\ne0$. H. Hopf conjectured that every even-dimensional compact
  manifold with positive sectional curvature has positive Euler characteristic.
  If the curvature operator is positive, this holds because then
  $\chi(M)=2$; if the Ricci curvature is positive, the fundamental group is
  finite, so $H_1(M,\mathbb{R})=0$, which already establishes the conjecture in
  dimension $2$. In dimension $4$, Poincar\'e duality with $H_1(M,\mathbb{R})=H_3(M,\mathbb{R})=0$
  gives $\chi(M)=1+\dim H_2(M,\mathbb{R})+1\ge2$.
\end{remark}

\begin{theorem}[Berger's zero theorem]
  \label{thm:pet-ch8-berger-killing-field-zero}
  \lean{PetersenLib.berger_evenDim_positiveCurvature_killingZero}
  \source{petersen:page-0335}
  \uses{def:pet-ch8-killing-field, prop:pet-ch8-killing-norm-formulas, prop:pet-ch8-killing-skew-symmetric}
  \dcref{ch8:8.3.1}
  If $(M,g)$ is a compact, even-dimensional manifold of positive sectional
  curvature, then every Killing field has a zero.
\end{theorem}
\begin{proof}
  \uses{thm:pet-ch8-berger-killing-field-zero}
  \sketch
  Suppose \(X\) is nowhere zero and minimize \(f=\tfrac12|X|^2\) at \(p\).
  Then \(\nabla f=-\nabla_XX=0\) by
  \cref{prop:pet-ch8-killing-norm-formulas}, so \(X_p\) belongs to the kernel of
  the skew-adjoint map \((\nabla X)_p\) from
  \cref{prop:pet-ch8-killing-skew-symmetric}. Because \(T_pM\) is
  even-dimensional, this kernel has even dimension and contains a vector
  \(v\) independent of \(X_p\). The Hessian formula gives
  \[
    \operatorname{Hess}f(v,v)=-R(v,X,X,v)<0,
  \]
  contradicting minimality.
\end{proof}

\begin{remark}[failure in odd dimensions]
  \label{rem:pet-ch8-berger-odd-dimension-counterexample}
  \lean{PetersenLib.berger_oddDim_counterexample}
  \source{petersen:page-0335}
  \uses{thm:pet-ch8-berger-killing-field-zero}
  \cref{thm:pet-ch8-berger-killing-field-zero} fails in odd dimensions: the
  unit vector field generating the Hopf fibration $S^3(1)\to S^2(1/2)$ is a
  Killing field on the positively curved, odd-dimensional manifold $S^3(1)$
  that is nowhere zero.
\end{remark}

\begin{definition}[Euler characteristic]
  \label{def:pet-ch8-euler-characteristic}
  \lean{PetersenLib.eulerCharacteristic}
  \source{petersen:page-0335}
  The \emph{Euler characteristic} of a space $M$ is the alternating sum
  \[
    \chi(M)=\sum_{p=0}^n(-1)^p\dim H_p(M,\mathbb{R}).
  \]
\end{definition}

\begin{theorem}[Euler characteristic from the Killing zero set]
  \label{thm:pet-ch8-euler-characteristic-zero-set-sum}
  \lean{PetersenLib.eulerChar_killingZeroSet_sum}
  \source{petersen:page-0335,petersen:page-0336}
  \uses{def:pet-ch8-killing-field, def:pet-ch8-euler-characteristic}
  \dcref{ch8:8.3.2}
  Let $X$ be a Killing field on a compact Riemannian manifold $M$, and let
  $N_i\subset M$ be the components of the zero set of $X$. Then
  $\chi(M)=\sum_i\chi(N_i)$.
\end{theorem}
\begin{proof}
  \uses{thm:pet-ch8-euler-characteristic-zero-set-sum}
  \sketch
  Choose disjoint tubular neighborhoods \(T_i\) of the zero components \(N_i\).
  The Killing flow preserves each distance to \(N_i\), so \(X\) is tangent to
  \(\partial T_i\) and is nonvanishing on the complement and on these
  boundaries. Poincar\'e--Hopf gives zero Euler characteristic for both pieces
  outside the tubes and their common boundary. Mayer--Vietoris therefore gives
  \[
    \chi(M)=\sum_i\chi(T_i)=\sum_i\chi(N_i),
  \]
  since every \(T_i\) deformation retracts onto \(N_i\).
\end{proof}

\begin{corollary}[positive Euler characteristic in dimension six]
  \label{cor:pet-ch8-positive-euler-characteristic-6manifold}
  \lean{PetersenLib.positiveCurvature_6manifold_killingField_eulerChar_pos}
  \source{petersen:page-0336}
  \uses{thm:pet-ch8-euler-characteristic-zero-set-sum, thm:pet-ch8-zero-set-totally-geodesic}
  \dcref{ch8:8.3.3}
  If $M$ is a compact $6$-manifold with positive sectional curvature that
  admits a Killing field, then $\chi(M)>0$.
\end{corollary}
\begin{proof}
  \uses{cor:pet-ch8-positive-euler-characteristic-6manifold}
  \sketch
  A zero exists by \cref{thm:pet-ch8-berger-killing-field-zero}. By
  \cref{thm:pet-ch8-zero-set-totally-geodesic}, every zero component is
  positively curved and has dimension \(0\), \(2\), or \(4\). Each has positive
  Euler characteristic, using the low-dimensional argument in
  \cref{rem:pet-ch8-hopf-conjecture-low-dimensions}. The formula in
  \cref{thm:pet-ch8-euler-characteristic-zero-set-sum} then makes
  \(\chi(M)\) a nonempty sum of positive integers.
\end{proof}

\begin{lemma}[orbit angle-sum bound for an isolated zero]
  \label{lem:pet-ch8-isolated-zero-orbit-angle-bound}
  \lean{PetersenLib.doublyWarpedOrbitAngleSumBound}
  \source{petersen:page-0336,petersen:page-0337}
  \uses{def:pet-ch8-killing-field}
  Let $p$ be an isolated zero of a Killing field $X$ on a Riemannian
  $4$-manifold, with flow $F^t$ inducing an isometric, fixed-point-free action
  $R^t=DF^t|_p$ on $S^3\subset T_pM$. Decomposing $T_pM=V_1\oplus V_2$ into the
  orthogonal invariant subspaces for this action exhibits $S^3$ as a doubly
  warped product over $[0,\pi/2]$, with associated distance function $r$
  measuring the angle to $V_1$; for all $v_1,v_2,v_3\in S^3$,
  \[
    \lvert r(v_1)r(v_2)\rvert+\lvert r(v_1)r(v_3)\rvert
    +\lvert r(v_2)r(v_3)\rvert\le\pi,
  \]
  and for every $\varepsilon>0$ there exist $t_1,t_2,t_3$ such that
  \[
    \angle(R^{t_1}v_1,R^{t_2}v_2)+\angle(R^{t_1}v_1,R^{t_3}v_3)+\angle(R^{t_2}v_2,R^{t_3}v_3)\le\pi+\varepsilon.
  \]
\end{lemma}
\begin{proof}
  \uses{lem:pet-ch8-isolated-zero-orbit-angle-bound}
  \sketch
  The linearized flow splits \(T_pM=V_1\oplus V_2\) into invariant
  \(2\)-planes and rotates them with speeds \(\theta_1,\theta_2\). The resulting
  doubly warped description of \(S^3\) gives the stated bound for the orbit
  distance function \(r\). If \(\theta_1/\theta_2\) is irrational, each orbit
  is dense in the corresponding product of level sets, so suitable flow times
  make the angle sum at most \(\pi+\varepsilon\). Rational speed ratios follow
  by approximation and continuity.
\end{proof}

\begin{lemma}[Toponogov angle-sum lower bound]
  \label{lem:pet-ch8-toponogov-angle-lower-bound}
  \lean{PetersenLib.toponogovTriangleAngleSumLowerBound}
  \source{petersen:page-0337}
  Let $M$ be a compact Riemannian $4$-manifold with $\sec\ge1$, and let
  $p_1,p_2,p_3,p_4\in M$ be four distinct points such that segments between any
  two of them exist. Writing $\inf\angle p\,x\,q$ for the infimum, over unit
  vectors in the directions of segments from $x$ to $p$ and from $x$ to $q$, of
  the angle between them, Toponogov's comparison theorem (comparing each
  triangle with $p_ip_jp_k$'s sides to a triangle in $S^2(1)$) gives, for each
  choice of three of the four points,
  \[
    \inf\angle p_ip_jp_k+\inf\angle p_kp_ip_j+\inf\angle p_jp_kp_i>\pi.
  \]
\end{lemma}
\begin{proof}
  \uses{lem:pet-ch8-toponogov-angle-lower-bound}
  \sketch
  Apply Toponogov comparison to each geodesic triangle determined by three of
  the four points. Every actual vertex angle is bounded below by the
  corresponding angle of its unit-sphere comparison triangle. Such a
  nondegenerate spherical triangle has positive spherical excess, so its three
  comparison angles sum to more than \(\pi\); taking the indicated infima
  preserves the lower bound.
\end{proof}

\begin{theorem}[Hsiang--Kleiner four-manifold theorem]
  \label{thm:pet-ch8-hsiang-kleiner-4manifold}
  \lean{PetersenLib.hsiangKleiner_4manifold_eulerChar_le_three}
  \source{petersen:page-0336}
  \uses{def:pet-ch8-killing-field, lem:pet-ch8-grove-searle-codim2, lem:pet-ch8-isolated-zero-orbit-angle-bound, lem:pet-ch8-toponogov-angle-lower-bound}
  \dcref{ch8:8.3.4}
  If $M^4$ is a compact, orientable, positively curved $4$-manifold that admits
  a Killing field, then $\chi(M)\le3$. In particular, $M$ is topologically
  equivalent to $S^4$ or $\mathbb{CP}^2$.
\end{theorem}
\begin{proof}
  \uses{thm:pet-ch8-hsiang-kleiner-4manifold}
  \sketch
  Rescale so that \(\sec\ge1\). A two-dimensional zero component is handled by
  \cref{lem:pet-ch8-grove-searle-codim2}. Otherwise all zeros are isolated.
  Four zeros would give, after summing at the four vertices, an angle total at
  most \(4\pi\) by
  \cref{lem:pet-ch8-isolated-zero-orbit-angle-bound}, but a total strictly
  greater than \(4\pi\) by
  \cref{lem:pet-ch8-toponogov-angle-lower-bound}. Hence there are at most three
  zeros. The formula
  \cref{thm:pet-ch8-euler-characteristic-zero-set-sum} gives
  \(\chi(M)\le3\), and the codimension-two classification yields the stated
  topological alternatives.
\end{proof}

\begin{theorem}[dependence of commuting fields in positive curvature]
  \label{thm:pet-ch8-grove-searle-commuting-dependent}
  \lean{PetersenLib.groveSearle_commutingKillingFields_dependentSomewhere}
  \source{petersen:page-0337,petersen:page-0338}
  \uses{def:pet-ch8-killing-field, prop:pet-ch8-killing-curvature-identities}
  \dcref{ch8:8.3.5}
  Let $M$ be a compact $n$-manifold with positive sectional curvature. Two
  commuting Killing fields must be linearly dependent somewhere on $M$.
\end{theorem}
\begin{proof}
  \uses{thm:pet-ch8-grove-searle-commuting-dependent}
  \sketch
  For commuting \(X,Y\), \cref{prop:pet-ch8-killing-curvature-identities}
  reduces the curvature identity to
  \[
    R(X,Y,Y,X)=-g(\nabla_XX,\nabla_YY)+|\nabla_XY|^2.
  \]
  If \(Y\) never vanishes, consider the squared norm of the component of \(X\)
  orthogonal to \(Y\),
  \[
    f=\frac12\left(|X|^2-\frac{g(X,Y)^2}{|Y|^2}\right).
  \]
  Were \(f>0\), at a minimum one may replace \(X\) by \(X-aY\) so that
  \(X\perp Y\). The norm Hessian formula from
  \cref{prop:pet-ch8-killing-norm-formulas}, the displayed curvature identity,
  and skew-adjointness of both covariant derivatives produce a direction in
  which \(\operatorname{Hess}f<0\), contradicting minimality. Thus \(f=0\)
  somewhere, exactly where \(X,Y\) are dependent.
\end{proof}

\begin{lemma}[orthogonal fields along a geodesic with drift towards the action]
  \label{lem:pet-ch8-jacobi-type-field}
  \lean{PetersenLib.wilkingOrthogonalField}
  \source{petersen:page-0338,petersen:page-0339}
  \uses{def:pet-ch8-killing-field}
  Let $M^n$ have positive sectional curvature, $X$ a Killing field, and let
  $N^{n-k}$ be a component of the zero set of $X$. Let $c:[0,1]\to M$ be a unit
  speed geodesic perpendicular to $N$ at both endpoints. Consider vector fields
  $E$ along $c$, orthogonal to both $c$ and the flow direction of $X$, defined
  by the linear ODE
  \[
    E(0)\in T_{c(0)}N,\qquad
    \dot E=-\frac{g(E,\nabla_{\dot c}X)}{|X|^2}X.
  \]
  Then $g(E,X)\equiv0$; $g(E(1),\nabla_{\dot c(1)}X)=0$; $g(E,\dot c)\equiv0$;
  and the space of such fields $E$ for which $E(1)\in T_{c(1)}N$ has dimension
  at least $n-2k+2$.
\end{lemma}
\begin{proof}
  \uses{lem:pet-ch8-jacobi-type-field}
  \sketch
  Skew-adjointness of \(\nabla X\) and the fact that
  \(T_pN=\ker(\nabla X)_p\), from
  \cref{thm:pet-ch8-zero-set-totally-geodesic}, make the endpoint singularities
  in the ODE removable. Differentiating the relevant inner products gives
  \[
    \frac d{dt}g(E,X)=0,\qquad
    \frac d{dt}g(X,\dot c)=0,\qquad
    \frac d{dt}g(E,\dot c)=0,
  \]
  which proves the three orthogonality assertions, including the limiting
  endpoint condition. Initial values range over an \((n-k)\)-space; imposing
  the endpoint normal constraints removes at most \(2(k-1)\) dimensions.
  Hence at least \(n-2k+2\) solutions end in \(T_{c(1)}N\).
\end{proof}

\begin{lemma}[Wilking's connectedness principle with symmetries]
  \label{lem:pet-ch8-wilking-connectedness-principle}
  \lean{PetersenLib.wilkingConnectednessPrincipleWithSymmetries}
  \source{petersen:page-0338,petersen:page-0339,petersen:page-0340}
  \uses{def:pet-ch8-killing-field, lem:pet-ch8-jacobi-type-field, thm:pet-ch8-zero-set-totally-geodesic}
  \dcref{ch8:8.3.6}
  Let $M^n$ be a compact $n$-manifold with positive sectional curvature and $X$
  a Killing field. If $N^{n-k}$ is a component of the zero set of $X$, then
  $N\subset M$ is $(n-2k+2)$-connected.
\end{lemma}
\begin{proof}
  \uses{lem:pet-ch8-wilking-connectedness-principle}
  \sketch
  By \cref{lem:pet-ch8-jacobi-type-field}, every geodesic perpendicular to
  \(N\) at both ends carries at least \(n-2k+2\) admissible fields \(E\).
  Shrink the metric in the \(X\)-orbit direction while fixing its orthogonal
  complement. The ODE makes the positive term in the index form tend to zero,
  whereas positive sectional curvature leaves
  \[
    -\int_0^1\sec(E,\dot c)|E|^2\,dt<0.
  \]
  Thus each such geodesic has index at least \(n-2k+2\). Morse theory for the
  path space relative to \(N\) converts this uniform index bound into
  \((n-2k+2)\)-connectivity of \(N\hookrightarrow M\).
\end{proof}

\begin{lemma}[codimension-two zero-set classification]
  \label{lem:pet-ch8-grove-searle-codim2}
  \lean{PetersenLib.groveSearle_codimTwoZeroSet_classification}
  \source{petersen:page-0340,petersen:page-0341}
  \uses{lem:pet-ch8-wilking-connectedness-principle, def:pet-ch8-euler-characteristic}
  \dcref{ch8:8.3.7}
  Let $M$ be a closed $n$-manifold with positive sectional curvature. If $M$
  admits a Killing field whose zero set has a component $N$ of codimension
  $2$, then $M$ is diffeomorphic to $S^n$, $\mathbb{CP}^{n/2}$, or a cyclic
  quotient of a sphere $S^n/\mathbb{Z}_q$.
\end{lemma}
\begin{proof}
  \uses{lem:pet-ch8-grove-searle-codim2}
  \sketch
  With codimension \(2\),
  \cref{lem:pet-ch8-wilking-connectedness-principle} makes
  \(N\hookrightarrow M\) \((n-2)\)-connected. For simply connected \(M\), the
  induced homology isomorphisms and Poincar\'e duality give two-periodicity:
  \[
    H_{j+2}(M)\simeq H_j(M),\qquad
    H^{j+2}(M)\simeq H^j(M)
  \]
  in the interior range. Odd \(n\) therefore gives the homology of a sphere;
  even \(n\) gives the cohomology of either a sphere or
  \(\mathbb{CP}^{n/2}\). The fixed-point tubular-neighborhood argument upgrades
  these alternatives to the asserted diffeomorphism types. On the universal
  cover the same analysis applies, and the deck group is cyclic, yielding the
  spherical quotients.
\end{proof}

\begin{proposition}[zero sets of commuting Killing fields]
  \label{prop:pet-ch8-commuting-killing-fields-properties}
  \lean{PetersenLib.commutingKillingFields_tangentToLevelSets,PetersenLib.commutingKillingFields_linearCombinationVanishesOnLarger}
  \source{petersen:page-0341}
  \uses{def:pet-ch8-killing-field, prop:pet-ch8-killing-curvature-identities}
  \dcref{ch8:8.3.8}
  Let $(M,g)$ be compact and $X,Y\in\iso(M,g)$ commute.
  \begin{enumerate}
    \item $Y$ is tangent to the level sets of $|X|^2$, in particular to the
      zero set of $X$.
    \item If $X$ and $Y$ both vanish on a totally geodesic submanifold
      $N\subset M$, then some linear combination of $X,Y$ vanishes on a
      strictly larger connected submanifold.
  \end{enumerate}
\end{proposition}
\begin{proof}
  \uses{prop:pet-ch8-commuting-killing-fields-properties}
  \sketch
  Since \([X,Y]=0\), invariance of both \(g\) and \(X\) under the \(Y\)-flow
  gives \(D_Y|X|^2=0\), proving the first assertion. At \(p\in N\),
  \cref{prop:pet-ch8-killing-curvature-identities} gives
  \([(\nabla X)_p,(\nabla Y)_p]=0\). These commuting skew-adjoint maps split
  the normal space of \(N\) into common invariant \(2\)-planes, while
  \cref{thm:pet-ch8-zero-set-totally-geodesic} identifies \(T_pN\) with their
  common kernel. On one normal block a nonzero linear combination of the two
  maps vanishes. The corresponding combination of \(X,Y\) therefore has a
  zero component whose tangent space strictly contains \(T_pN\).
\end{proof}

\begin{definition}[symmetry rank]
  \label{def:pet-ch8-symmetry-rank}
  \lean{PetersenLib.symmetryRank}
  \source{petersen:page-0341}
  \uses{def:pet-ch8-killing-field}
  The \emph{symmetry rank} of a Riemannian manifold $(M,g)$ is the dimension of
  a maximal subspace (equivalently, Abelian subalgebra) of commuting Killing
  fields in $\iso(M,g)$.
\end{definition}

\begin{theorem}[Grove--Searle maximal symmetry-rank classification]
  \label{thm:pet-ch8-symmetry-rank-half-classification}
  \lean{PetersenLib.groveSearle_symmetryRank_atLeastHalf_classification}
  \source{petersen:page-0341,petersen:page-0342}
  \uses{def:pet-ch8-symmetry-rank, lem:pet-ch8-grove-searle-codim2, prop:pet-ch8-commuting-killing-fields-properties}
  \dcref{ch8:8.3.9}
  Let $M$ be a compact $n$-manifold with positive sectional curvature and
  symmetry rank $k$. If $k\ge n/2$, then $M$ is diffeomorphic to a sphere, a
  complex projective space, or a cyclic quotient of a sphere $S^n/\mathbb{Z}_q$,
  with $\mathbb{Z}_q$ acting by isometries on the unit sphere.
\end{theorem}
\begin{proof}
  \uses{thm:pet-ch8-symmetry-rank-half-classification}
  \sketch
  Let \(\mathfrak a\) be an Abelian algebra of dimension \(k\ge n/2\).
  Repeated use of
  \cref{prop:pet-ch8-commuting-killing-fields-properties} produces a maximal
  nontrivial zero component \(N\) with
  \(\dim(\mathfrak a|_N)=k-1\). If its codimension exceeded \(2\), iteration
  inside zero components would end in dimension \(1\) or \(2\) with at least
  two independent restricted fields. This contradicts the isometry-algebra
  bound in dimension \(1\), or first-jet uniqueness from
  \cref{prop:pet-ch8-killing-uniqueness} in dimension \(2\). Hence
  \(\operatorname{codim}N=2\), and
  \cref{lem:pet-ch8-grove-searle-codim2} gives the classification.
\end{proof}

\begin{theorem}[Euler characteristic under intermediate symmetry rank]
  \label{thm:pet-ch8-puttmann-searle-rong}
  \lean{PetersenLib.puttmannSearleRong_symmetryRank_eulerChar_pos}
  \source{petersen:page-0342}
  \uses{def:pet-ch8-symmetry-rank, lem:pet-ch8-grove-searle-codim2, cor:pet-ch8-positive-euler-characteristic-6manifold}
  \dcref{ch8:8.3.10}
  If $M^{2n}$ is a compact $2n$-manifold with positive sectional curvature and
  symmetry rank $k\ge n/2-1$, then $\chi(M)>0$.
\end{theorem}
\begin{proof}
  \uses{thm:pet-ch8-puttmann-searle-rong}
  \sketch
  The dimensions \(2,4\) follow from
  \cref{rem:pet-ch8-hopf-conjecture-low-dimensions}, dimension \(6\) from
  \cref{cor:pet-ch8-positive-euler-characteristic-6manifold}, and the
  codimension-two case from \cref{lem:pet-ch8-grove-searle-codim2}. For the
  induction, strengthen the claim to positivity of the Euler characteristic
  of every zero-set component. If all maximal components have codimension at
  least \(4\), restriction lowers the Abelian rank by at most one and leaves
  enough rank to apply the strengthened hypothesis in lower dimension. If a
  codimension-two component occurs, the codimension-two lemma forces
  vanishing of odd Betti numbers for it and for \(M\). In either case every
  zero component has positive Euler characteristic, and summing them proves
  \(\chi(M)>0\).
\end{proof}

\begin{remark}[a positively curved exception with lower symmetry rank]
  \label{rem:pet-ch8-f4-spin8-exception}
  \lean{PetersenLib.f4Spin8SymmetryRankException}
  \source{petersen:page-0342}
  \uses{def:pet-ch8-symmetry-rank}
  \cref{thm:pet-ch8-puttmann-searle-rong} does not cover all known examples:
  the homogeneous space $F_4/\mathrm{Spin}(8)$ is a positively curved
  $24$-manifold with symmetry rank $4$.
\end{remark}

\begin{theorem}[Wilking's high symmetry-rank classification]
  \label{thm:pet-ch8-wilking-symmetry-rank-classification}
  \lean{PetersenLib.wilking_symmetryRank_classification}
  \source{petersen:page-0343}
  \uses{def:pet-ch8-symmetry-rank}
  \dcref{ch8:8.3.11}
  Let $M$ be a compact, simply connected, positively curved $n$-manifold with
  symmetry rank $k$, and $n\ge10$. If $k\ge n/4+1$, then $M$ has the topology
  of a sphere, complex projective space, or quaternionic projective space.
  Moreover, when $M$ is not simply connected, its fundamental group is cyclic.
\end{theorem}

\begin{theorem}[Kennard's Euler-characteristic criterion]
  \label{thm:pet-ch8-kennard-2013-euler-characteristic}
  \lean{PetersenLib.kennard2013_symmetryRank_eulerChar_pos}
  \source{petersen:page-0343}
  \uses{def:pet-ch8-symmetry-rank}
  \dcref{ch8:8.3.12}
  If $M^{4n}$ is a compact $4n$-manifold with positive sectional curvature and
  symmetry rank $k>2\log_4(4n)-2$, then $\chi(M)>0$.
\end{theorem}

\begin{theorem}[Betti-number bounds under large symmetry rank]
  \label{thm:pet-ch8-kennard-amann-betti-numbers}
  \lean{PetersenLib.kennardAmannKennard_bettiNumberBounds}
  \source{petersen:page-0343}
  \uses{def:pet-ch8-symmetry-rank}
  \dcref{ch8:8.3.13}
  Let $M^{2n}$ be a compact positively curved manifold with symmetry rank
  $k\ge\log_{4/3}(2n)$.
  \begin{enumerate}
    \item The Betti numbers satisfy $b_2\le1$ and $b_4\le1$.
    \item If $b_4=0$, then $\chi(M)=2$.
    \item Consequently, $M$ cannot have the homotopy type of a product
      $N\times N$ for compact simply connected $N$.
  \end{enumerate}
\end{theorem}

\section{Exercises}

The exercises develop completeness, explicit symmetry algebras, integral
identities, symmetry rank, and covering-space descriptions of isometries.

\begin{remark}[incomplete metrics and isometry algebras]
  \label{rem:pet-ch8-ex-8-4-1}
  \lean{PetersenLib.exercise_8_4_1}
  \source{petersen:page-0343}
  \uses{thm:pet-ch8-isometry-lie-algebra-dimension}
  \dcref{ch8:8.4.1}
  Show that \cref{thm:pet-ch8-isometry-lie-algebra-dimension} does not
  necessarily extend to incomplete Riemannian manifolds.
\end{remark}

\begin{remark}[parallel first derivative along a zero component]
  \label{rem:pet-ch8-ex-8-4-2}
  \lean{PetersenLib.exercise_8_4_2}
  \source{petersen:page-0343}
  \uses{def:pet-ch8-killing-field, thm:pet-ch8-zero-set-totally-geodesic}
  \dcref{ch8:8.4.2}
  Let $N$ be a component of the zero set of a Killing field $X$. Show that
  $\nabla_V(\nabla X)=0$ for vector fields $V$ tangent to $N$.
\end{remark}

\begin{remark}[coordinate Killing fields]
  \label{rem:pet-ch8-ex-8-4-3}
  \lean{PetersenLib.exercise_8_4_3}
  \source{petersen:page-0343}
  \uses{def:pet-ch8-killing-field}
  \dcref{ch8:8.4.3}
  Show that a coordinate vector field $\partial_i$ is a Killing field if and
  only if $\partial_kg_{ij}=0$.
\end{remark}

\begin{remark}[restriction and projection to submanifolds]
  \label{rem:pet-ch8-ex-8-4-4}
  \lean{PetersenLib.exercise_8_4_4}
  \source{petersen:page-0343}
  \uses{def:pet-ch8-killing-field}
  \dcref{ch8:8.4.4}
  Let $X$ be a Killing field on $(M,g)$ and $N\subset M$ a submanifold with the
  induced metric.
  \begin{enumerate}
    \item Show that if $X$ is tangent to $N$, then $X|_N$ is a Killing field
      on $N$.
    \item Show that if $N$ is totally geodesic, then $X^\top$, the projection
      of $X$ onto $N$, is a Killing field.
  \end{enumerate}
\end{remark}

\begin{remark}[completeness of Killing fields]
  \label{rem:pet-ch8-ex-8-4-5}
  \lean{PetersenLib.exercise_8_4_5}
  \source{petersen:page-0343,petersen:page-0344}
  \uses{def:pet-ch8-killing-field}
  \dcref{ch8:8.4.5}
  Let $(M,g)$ be a complete Riemannian manifold and $X$ a Killing field on $M$.
  \begin{enumerate}
    \item Let $c:[a,b]\to M$ be a geodesic. Show that there is $\varepsilon>0$
      and a geodesic variation $\bar c:(-\varepsilon,\varepsilon)\times[a,b]\to M$
      such that the curves $s\mapsto\bar c(s,t)$ are integral curves of $X$.
    \item Use completeness to extend $\bar c$ to
      $(-\varepsilon,\varepsilon)\times\mathbb{R}\to M$, and show that for all
      $t\in\mathbb{R}$ the curves $s\mapsto\bar c(s,t)$ are integral curves of
      $X$.
    \item Show that the integral curve of $X$ through any point of $M$ exists
      on $(-\varepsilon,\varepsilon)$.
    \item Show that $X$ is complete.
  \end{enumerate}
\end{remark}

\begin{remark}[maximal isometry dimension under nonpositive Ricci curvature]
  \label{rem:pet-ch8-ex-8-4-6}
  \lean{PetersenLib.exercise_8_4_6}
  \source{petersen:page-0344}
  \uses{def:pet-ch8-symmetry-rank, thm:pet-ch8-bochner-negative-ricci}
  \dcref{ch8:8.4.6}
  Let $(M^n,g)$ be a compact Riemannian $n$-manifold with $\dim(\iso(M,g))=n$
  and $\mathrm{Ric}\le0$. Show that $M$ is flat and $\pi_1=\mathbb{Z}^n$. (Hint:
  show $\widetilde M=\mathbb{R}^n$ and that the deck transformations have
  differential $I$, since they preserve a parallel orthonormal frame.)
\end{remark}

\begin{remark}[Killing algebras of the model spaces]
  \label{rem:pet-ch8-ex-8-4-7}
  \lean{PetersenLib.exercise_8_4_7}
  \source{petersen:page-0344}
  \uses{def:pet-ch8-killing-field}
  \dcref{ch8:8.4.7}
  Let $x^i$ be standard Cartesian coordinates on $\mathbb{R}^n$, and let
  $W=\mathrm{span}_{\mathbb{R}}\{1,x^1,\dots,x^n\}$,
  $V=\mathrm{span}_{\mathbb{R}}\{x^1,\dots,x^n\}$.
  \begin{enumerate}
    \item Show that $\iso(\mathbb{R}^n)$ is naturally isomorphic to $\Lambda^2W$,
      identifying $u\wedge v$ with the vector field $u\nabla v-v\nabla u$.
    \item Show similarly that $\iso(S^{n-1}(R))$ is naturally isomorphic to
      $\Lambda^2V$.
    \item Using $\nabla u=g^{ij}\partial_iu\,\partial_j$ for a
      pseudo-Riemannian space, redo (1) for $\mathbb{R}^{p,q}$ and (2) for
      $H^{n-1}(R)\subset\mathbb{R}^{n,1}$.
  \end{enumerate}
\end{remark}

\begin{remark}[conformal fields on the sphere]
  \label{rem:pet-ch8-ex-8-4-8}
  \lean{PetersenLib.exercise_8_4_8}
  \source{petersen:page-0344}
  \uses{def:pet-ch8-killing-field}
  \dcref{ch8:8.4.8}
  Let $x^i$ be standard Cartesian coordinates on $\mathbb{R}^{n+1}$ and
  $V=\mathrm{span}_{\mathbb{R}}\{1,x^1,\dots,x^{n+1}\}$, restricted to $S^n$.
  Show that for $u,v\in V$ the fields $u\nabla v-v\nabla u$ are conformal (a
  field $X$ is conformal if $L_Xg=\lambda g$ for some function $\lambda$).
\end{remark}

\begin{remark}[distribution generated by Killing fields]
  \label{rem:pet-ch8-ex-8-4-9}
  \lean{PetersenLib.exercise_8_4_9}
  \source{petersen:page-0344}
  \uses{prop:pet-ch8-killing-uniqueness}
  \dcref{ch8:8.4.9}
  Let $(M,g)$ be a Riemannian manifold and $t_p=\{X|_p\in T_pM\mid X\in\iso\}$.
  \begin{enumerate}
    \item Show that $t_p$ defines an integrable distribution on an open set
      $O\subset M$. (Hint: show
      $\{p\in M\mid\dim t_p=\max_{q\in M}\dim t_q\}$ is open.)
    \item Give an example where $O\ne M$.
    \item Show that $DF(t_p)=t_{F(p)}$ for all $F\in\mathrm{Iso}(M,g)$.
    \item Assume $(M,g)$ is complete. Show that the leaves of this
      distribution are homogeneous and properly embedded. (Hint: show the
      connected subgroup $\mathrm{Iso}_0\subset\mathrm{Iso}(M,g)$ containing
      the identity is closed and preserves the leaves.)
    \item Show that if $(M,g)$ is homogeneous, then $t_p=T_pM$ for all
      $p\in M$.
  \end{enumerate}
\end{remark}

\begin{remark}[rational lattice of a torus action]
  \label{rem:pet-ch8-ex-8-4-10}
  \lean{PetersenLib.exercise_8_4_10}
  \source{petersen:page-0344}
  \uses{def:pet-ch8-killing-field}
  \dcref{ch8:8.4.10}
  Let $\mathfrak t\subset\iso(M,g)$ be an Abelian subalgebra corresponding to a
  torus subgroup $T^k\subset\mathrm{Iso}(M,g)$. Define $\mathfrak p\subset\mathfrak t$
  as the Killing fields corresponding to circle actions, i.e.\ actions induced
  by homomorphisms $S^1\to T^k$. Show that $\mathfrak p$ is a vector space over
  $\mathbb{Q}$ with $\dim_{\mathbb{Q}}\mathfrak p=\dim_{\mathbb{R}}\mathfrak t$.
\end{remark}

\begin{remark}[Laplacian of an inner product of Killing fields]
  \label{rem:pet-ch8-ex-8-4-11}
  \lean{PetersenLib.exercise_8_4_11}
  \source{petersen:page-0344}
  \uses{def:pet-ch8-killing-field}
  \dcref{ch8:8.4.11}
  Given two Killing fields $X,Y$ on a Riemannian manifold, develop a formula
  for $\Delta g(X,Y)$. Use this to give a formula for the Ricci curvature in a
  frame consisting of Killing fields.
\end{remark}

\begin{remark}[integral characterizations of Killing fields]
  \label{rem:pet-ch8-ex-8-4-12}
  \lean{PetersenLib.exercise_8_4_12}
  \source{petersen:page-0345}
  \uses{def:pet-ch8-killing-field}
  \dcref{ch8:8.4.12}
  Let $X$ be a vector field on a Riemannian manifold.
  \begin{enumerate}
    \item Show that
      $|L_Xg|^2=2|\nabla X|^2+2\operatorname{tr}(\nabla X\circ\nabla X)$.
    \item Establish, on a closed oriented Riemannian manifold,
      \[
        \int_M\Big(\mathrm{Ric}(X,X)+\operatorname{tr}(\nabla X\circ\nabla X)-(\operatorname{div}X)^2\Big)=0,
      \]
      \[
        \int_M\Big(\mathrm{Ric}(X,X)+g(\operatorname{tr}\nabla^2X,X)+\tfrac12|L_Xg|^2-(\operatorname{div}X)^2\Big)=0.
      \]
    \item Show that $X$ is a Killing field if and only if
      $\operatorname{div}X=0$ and $\operatorname{tr}\nabla^2X=-\mathrm{Ric}(X)$.
  \end{enumerate}
\end{remark}

\begin{remark}[Yano's theorem for affine vector fields]
  \label{rem:pet-ch8-ex-8-4-13}
  \lean{PetersenLib.exercise_8_4_13}
  \source{petersen:page-0345}
  \uses{rem:pet-ch8-ex-8-4-12}
  \dcref{ch8:8.4.13}
  If $X$ is an affine vector field, show that
  $\operatorname{tr}\nabla^2X=-\mathrm{Ric}(X)$ and that $\operatorname{div}X$
  is constant. Use this together with the characterizations of
  \cref{rem:pet-ch8-ex-8-4-12} to show that on closed manifolds affine fields
  are Killing fields.
\end{remark}

\begin{remark}[commutation with the scalar Laplacian]
  \label{rem:pet-ch8-ex-8-4-14}
  \lean{PetersenLib.exercise_8_4_14}
  \source{petersen:page-0345}
  \uses{def:pet-ch8-killing-field}
  \dcref{ch8:8.4.14}
  Let $X$ be a vector field on a Riemannian manifold. Show that $X$ is a
  Killing field if and only if $L_X$ and $\Delta$ commute on functions.
\end{remark}

\begin{remark}[basic symmetry-rank bounds and examples]
  \label{rem:pet-ch8-ex-8-4-15}
  \lean{PetersenLib.exercise_8_4_15}
  \source{petersen:page-0345}
  \uses{def:pet-ch8-symmetry-rank}
  \dcref{ch8:8.4.15}
  Let $(M,g)$ be a compact $n$-manifold with positive sectional curvature and
  $\mathfrak a\subset\iso$ an Abelian subalgebra.
  \begin{enumerate}
    \item Show that $\dim\mathfrak a\le n/2$.
    \item Show that spheres and complex projective spaces have maximal
      symmetry rank.
    \item Show that the flat torus $T^n$ has symmetry rank $n$.
  \end{enumerate}
\end{remark}

\begin{remark}[maximal zero components for commuting fields]
  \label{rem:pet-ch8-ex-8-4-16}
  \lean{PetersenLib.exercise_8_4_16}
  \source{petersen:page-0345}
  \uses{def:pet-ch8-symmetry-rank, prop:pet-ch8-commuting-killing-fields-properties}
  \dcref{ch8:8.4.16}
  Let $(M,g)$ be a compact $n$-manifold with positive sectional curvature and
  $\mathfrak a\subset\iso$ an Abelian subalgebra. Let $\mathcal Z(\mathfrak a)$
  be the set of nontrivial connected components of zero sets of Killing fields
  in $\mathfrak a$. Show that if $N\in\mathcal Z(\mathfrak a)$ is maximal under
  inclusion, then $\dim(\mathfrak a|_N)=\dim\mathfrak a-1$ and
  $\dim N\ge2(\dim\mathfrak a-1)$.
\end{remark}

\begin{remark}[Kennard's high-rank homology argument]
  \label{rem:pet-ch8-ex-8-4-17}
  \lean{PetersenLib.exercise_8_4_17}
  \source{petersen:page-0345,petersen:page-0346}
  \uses{def:pet-ch8-symmetry-rank, thm:pet-ch8-symmetry-rank-half-classification, rem:pet-ch8-ex-8-4-16, lem:pet-ch8-wilking-connectedness-principle}
  \dcref{ch8:8.4.17}
  Let $(M^n,\rho)$ be a simply connected compact $n$-manifold with positive
  sectional curvature and symmetry rank $\ge3n/8+1$.
  \begin{enumerate}
    \item When $n\le8$, conclude that the assumptions are covered by
      \cref{thm:pet-ch8-symmetry-rank-half-classification}.
    \item Use \cref{rem:pet-ch8-ex-8-4-16} to show that if
      $N\in\mathcal Z(\mathfrak a)$ is maximal under inclusion, then
      $\dim N\ge3n/4$.
    \item Show that a maximal $N\in\mathcal Z(\mathfrak a)$ either has
      codimension $2$ or has symmetry rank $\ge\tfrac{3\dim N}8+1$.
    \item Use induction on $n$ to show that $M$ has the homology groups of a
      sphere or complex projective space. (Hint: when the maximal
      $N\in\mathcal Z(\mathfrak a)$ has codimension $\ge4$, use
      \cref{lem:pet-ch8-wilking-connectedness-principle} to show $N$ is simply
      connected, then use it again to compute homology/homotopy groups in
      dimensions $\le n/2$; finally use Poincar\'e duality to find all
      homology groups of $M$.)
  \end{enumerate}
\end{remark}

\begin{remark}[isometries through a universal covering]
  \label{rem:pet-ch8-ex-8-4-18}
  \lean{PetersenLib.exercise_8_4_18}
  \source{petersen:page-0346}
  \uses{def:pet-ch8-killing-field}
  \dcref{ch8:8.4.18}
  Let $\widetilde M\to M$ be the universal covering, with deck transformations
  $\Gamma=\pi_1(M)$ acting as isometries on $\widetilde M$.
  \begin{enumerate}
    \item Show that
      $\iso(M)=\{X\in\iso(\widetilde M)\mid G_*X=X\text{ for all }G\in\Gamma\}$.
    \item Show that the connected Lie subgroup corresponding to
      $\iso(M)\subset\iso(\widetilde M)$ is the identity component of the
      centralizer
      $\mathbf C(\Gamma)=\{F\in\mathrm{Iso}(\widetilde M)\mid FG=GF\text{ for all }G\in\Gamma\}$.
    \item Show that $\mathrm{Iso}(M)$ can be identified with
      $\mathbf N(\Gamma)/\Gamma$, where
      $\mathbf N(\Gamma)=\{F\in\mathrm{Iso}(\widetilde M)\mid F\Gamma F^{-1}=\Gamma\}$
      is the normalizer.
  \end{enumerate}
\end{remark}
